%!TEX root = ../main.tex
\section{Результаты и их обсуждение}

Чтобы исключить фактор случайности, я прогнал все 6 методов на 20 независимых случайных сидах (10 эпизодов на задачу, 5 задач, суммарно 1000 эпизодов на метод). В лабиринте Medium успешными считались эпизоды длительностью менее 1000 шагов, а в Large --- 1500. Результаты сведены в табл.~\ref{tab:main_benchmark}.


\begin{table*}[t]
    \centering
    \small
    \setlength{\tabcolsep}{8pt}
    \caption{Сводные результаты всех 6 методов на AntMaze Medium и Large (20 сидов, 1000 эпизодов)}
    \label{tab:main_benchmark}
    \begin{tabular}{lcccc}
        \toprule
        \multirow{2}{*}{\textbf{Метод планирования}} & \multicolumn{2}{c}{\textbf{AntMaze Medium}} & \multicolumn{2}{c}{\textbf{AntMaze Large}} \\
        \cmidrule(lr){2-3} \cmidrule(lr){4-5}
        & \textbf{SR (\%)} & \textbf{Время (мс)} & \textbf{SR (\%)} & \textbf{Время (мс)} \\
        \midrule
        1. Single-Intention Baseline & $81.4 \pm 4.5$ & 0.95 & $49.3 \pm 5.2$ & 0.95 \\
        2. Recursive Bisection & $75.9 \pm 6.4$ & 2.82 & $46.5 \pm 8.0$ & 2.36 \\
        3. Dijkstra Teacher (\texttt{high\_actor}) & $83.9 \pm 3.2$ & 0.85 & $48.7 \pm 9.2$ & 0.87 \\
        4. Distilled JAX GatedAttn [$O(1)$] & $\mathbf{84.3 \pm 4.8}$ & 0.55 & $45.6 \pm 5.8$ & 0.57 \\
        5. Dijkstra + Single WP Translator & $83.3 \pm 5.0$ & \textbf{0.53} & $57.1 \pm 5.6$ & \textbf{0.54} \\
        \textbf{6. Dijkstra + Enhanced Sequence Attn} & $83.6 \pm 3.8$ & 0.95 & $\mathbf{60.9 \pm 6.9}$ & 0.97 \\
        \bottomrule
    \end{tabular}
\end{table*}

\subsection{Лабиринт Medium}
Лабиринт Medium относительно компактный, коридоры короткие. Здесь почти все алгоритмы отработали одинаково хорошо. Исключением стал алгоритм бисекции: он показал результат хуже всех, отстав почти на $10\%$. Все алгоритмы на основе Search on the Replay Buffer показали себя лучше бейзлайна на $2\text{--}3\%$.

\subsection{Лабиринт Large}
Лабиринт Large в три раза больше по линейным размерам, чем Medium, и поворотов здесь значительно больше. Методы на основе Search on the Replay Buffer показали себя ещё лучше, а в комбинации со специализированными трансляторами уверенно превзошли бейзлайн: одноточечный транслятор — на $8\%$, а трансформер над последовательностью вейпоинтов — на $11\%$. То, что дистиллированная версия алгоритма Дейкстры показала себя слабее на Large, объясняется тем, что на большом лабиринте критически необходимо явно планировать маршрут глубже, чем на одну точку вперёд, чего реактивная сеть без графа не делает.

\subsection{Итог}
На лабиринте Medium преимущество над бейзлайном составило около $3\%$, так как лабиринт недостаточно велик, чтобы выигрыш от долгосрочного планирования проявился в полной мере. Чего нельзя сказать про Large, где методы с явным планированием и учётом геометрии траектории вырвались вперёд на $7\text{--}11\%$.